\documentclass[10pt,a4paper]{amsart}
\usepackage[margin=19mm]{geometry}
\usepackage{amsmath,amssymb,mathtools}
\usepackage[scr=rsfs,cal=euler]{mathalfa}
\usepackage{xfrac}
\usepackage[unicode,hidelinks,bookmarks=false]{hyperref}
\newcommand{\ie}{i.e.\ }
\newcommand{\cf}{cf.\ }
\newcommand{\resp}{respectively\ }
\newcommand{\Subsection}{\subsection}
\providecommand{\nobreakdash}{\nobreak\hbox{-}}
\setlength{\emergencystretch}{2em}
\multlinegap0pt
\usepackage{bm}
\usepackage{bbm}
\usepackage{dsfont}
\usepackage{xspace}
\usepackage{cancel}
\usepackage{mathtools}
\usepackage{amsthm}
\makeatletter
\@ifundefined{lemma}{\newtheorem{lemma}{Lemma}}{}
\makeatother
\newcommand{\R}{\mathbb{R}}
\newcommand{\Z}{\mathbb{Z}}
\newcommand{\eps}{\varepsilon}
\title[A method of solution for the inverse problem for $h$-functio]{A method of solution for the inverse problem for $h$-functions of planar Brownian motion}
\author{Greg Markowsky, Clayton McDonald}
\date{Source: arXiv:2412.05764v2 (CC BY 4.0)}
\begin{document}
\maketitle

\noindent\emph{Source and licence.} A method of solution for the inverse problem for $h$-functions of planar Brownian motion. Version arXiv:2412.05764v2 (CC BY 4.0), distributed under the Creative Commons Attribution 4.0 International licence, as recorded in the arXiv licence metadata for this exact version and shown on the abstract page https://arxiv.org/abs/2412.05764v2. Full rights notice: licenses/arxiv-2412.05764-CC-BY-4.0.txt.\par\smallskip

\noindent\textit{Editorial scope.} Counted unit: the Lemma whose source label is \texttt{lemma:2} in the article (source0.tex lines 331--338, source label \texttt{lemma:2}), delivered together with the article's own complete original proof of it (source0.tex lines 339--382, one \texttt{proof} environment of 44 lines). No mathematics is written, repaired or re-derived: statement and proof are the article's own text line for line. Required context retained uncounted: none. Every definition, display and label of those lines is retained unchanged. Omitted from this package and not redistributed: 1--330, 383--745. Nothing inside a retained excerpt is omitted or altered: every retained body is delivered exactly as the article's lines are written, so no omission receipt of any kind accompanies this package. Citations: the retained excerpts carry no citation command at all, so no bibliography entry of the article is transcribed and no reference is redistributed. Reference adaptation: none was needed and none was made; every reference inside the delivered unit resolves inside it, so no unresolved reference or citation is produced. Printed slips kept literal: no printed word, symbol, hypothesis or number of the counted statement or of its proof was altered. Layout and class: the article's own class is \texttt{amsart} with packages including geometry, hyperref, amsfonts, amsmath, amssymb, amsthm, float, tikz-cd, svg, natbib; the wrapper is the lane's pinned \texttt{amsart} editorial wrapper at 10pt on A4, which restates the theorem-like environments the retained text needs and adds the article's own macro definitions that the retained text needs (\texttt{\textbackslash R}, \texttt{\textbackslash Z}, \texttt{\textbackslash eps}), with extra wrapper packages (bm, bbm, dsfont, xspace, cancel, mathtools). The delivered numbering is the unit's own. Assets: no retained span contains a figure environment, a graphics-inclusion command, a PDF-page inclusion, a table or a TikZ picture; no image file is copied into this package, the package declares no asset dependency and no image file is redistributed with it. Licence: version arXiv:2412.05764v2 of this article is distributed under the Creative Commons Attribution 4.0 International licence, as recorded in the arXiv licence metadata for this exact version and shown on the abstract page https://arxiv.org/abs/2412.05764v2. A full rights notice is delivered at licenses/arxiv-2412.05764-CC-BY-4.0.txt. Characters: every retained excerpt is pure ASCII, so the delivered engine, XeLaTeX with its default Latin Modern text font, reports no missing character.
\begin{lemma}\label{lemma:2}
    The periodic Hilbert transform takes the form
    \begin{align*}
        (Hf)(x) = \frac{1}{2} \mathrm{p.v.}\int_T f(t) \cot \frac{\pi}{2} (x-t)dt.
    \end{align*}
    
    This exists for $f\in L^p(T)$, and furthermore $H$ is bounded on $L^p(T)$, for $1<p<\infty$, and commutes with the Poisson integral. 
\end{lemma}

\begin{proof}
 First we prove existence. Let $f\in L^p(T)$. Then
 
    \begin{align*}
        (H f)(x) = \frac{1}{\pi} \mathrm{p.v.} \int_\R \frac{f(t)}{x-t} dt =& \frac{1}{\pi} 
	    \lim_{\eps\to 0}\int_{|t| > \eps} \frac{f(x-t)}{t}dt \\                                         
	    =& \frac{1}{\pi} \lim_{\eps\to 0} \int_{\eps \leq |t| \leq 1} \frac{f(x-t)}{t}dt +
	\frac{1}{\pi}\int_{-1}^1 \sum_{\substack{k\in\Z\\k\neq0}}\frac{f(x-t)}{t-2k}dt 
    \end{align*}
    This follows by the periodicity of $f$. The Hilbert transform of an
    $L^p(\R)$ function exists almost everywhere. Let $f_1 :=
    \mathbf{1}_{[-2, 2]} f$, so $f_1\in L^p(\R)$ and on $(-2, 2)$ we have
    $f_1 = f$. That is, the first integral in the last equation is
    nothing but the Hilbert transform of an $L^p(\R)$ function, hence it exists
    for almost every $x$. 
    Again a routine application of the residue theorem yields the following equality.
    \begin{align*}
        \sum_{\substack{k\in\Z\\k\neq0}} \frac{1}{t-2k} = \frac{\pi}{2}
	    \cot \frac{\pi}{2}t - \frac{1}{t}
    \end{align*}
    This is bounded on $(-1, 1)$, and so the second integral term is finite for
    almost all $x$. This proves existence, the form trivially follows by
    replacing $\int_{-1}^1 f(x-t) \left(\frac{\pi}{2}\cot \frac{\pi}{2}t
    - \frac{1}{t}\right)dt$ with 
    $\lim_{\eps\to 0} \int_{\eps \leq |t| \leq 1}  f(x-t) 
	\left(\frac{\pi}{2}\cot \frac{\pi}{2}t - \frac{1}{t}\right)dt$ and simplifying.\\
    
    Now, to prove boundedness on $L^p(T)$ for $1 < p < \infty$. Let $H_\eps f :=  
	\frac{1}{\pi} \int_{|t| > \eps} \frac{ f(x-t)}{t}dt$, then
    \begin{align*}
        \Vert H_\eps  f\Vert_p &\leq \bigg\Vert \frac{1}{\pi} \int_{\eps \leq |t| \leq 1} 
	    \frac{ f(x-t)}{t}dt \bigg \Vert_p + 
        \bigg\Vert \frac{1}{\pi}\int_{-1}^1 \sum_{\substack{k\in\Z\\k\neq0}}
	    \frac{ f(x-t)}{t-2k}dt\bigg\Vert_p  \\ &\leq A\Vert f_1\Vert_p + B \Vert  
	    f\Vert_p \leq C\Vert f\Vert_p 
    \end{align*}
    The first bound follows as the Hilbert transform is bounded on $\R$ and $f_1\in L^p(\R)$ agrees 
	with $f$ on $(-1, 1)$. The second term follows by repeated application of the triangle 
	inequality. Fatou's lemma with $\eps \to 0$ gives boundedness on $L^p$ 
	for $1 < p < \infty$. \\

    Finally, as the Hilbert transform is translationally invariant it must commute with 
	convolution operators, so for $p\in(1, \infty)$ $P_y\ast Hf = H (P_y \ast f)$.
\end{proof}

\end{document}
